\chapter{Guide to the Worked-Example Code}
\label{app:e}

Each chapter has a corresponding Python file, located at \texttt{code/chapter\_XX.py}. These scripts generate the PDF figures in the main text and also provide a starting point for reproducing the order-of-magnitude estimates of each experiment. Formula definitions follow their first appearance in the main text; the code turns those formulas into checkable curves, scalings, simulated data, and error budgets.

\section{Contents and Run Order}
\label{appsec:e-run}

\begin{enumerate}[leftmargin=*]
\item First read \texttt{code/README.md} to confirm the Python version, the required packages, and the output directory.
\item Look at \texttt{code/plot\_style.py} and \texttt{code/plot\_recipes.py} to understand the unified fonts, colors, line widths, and save functions.
\item Run a single-chapter script, for example \texttt{python code/chapter\_20.py}. A single-chapter script should generate only that chapter's figures and should not modify other chapters.
\item Check that the PDFs are written to \texttt{figures/generated/chapter\_XX/}.
\item Compile \texttt{main\_nature.tex} and confirm that the figure numbers, captions, in-text references, and PDF layout are all consistent.
\end{enumerate}

\section{Correspondence Between Scripts and Text}
\label{appsec:e-map}

\begin{longtable}{p{0.18\linewidth}p{0.30\linewidth}p{0.42\linewidth}}
\toprule
Chapter & Typical figure or computation & Points to check \\
\midrule
1--2 & Event tables, Poisson counts, ordinary light curves, and information loss & Event-table fields in the figures match the symbols in the text. \\
3--4 & Light states, occupation numbers, \(g^{(1)}\), \(g^{(2)}\), multimode dilution & Do not draw idealized single-mode results as if they were real instrumental observations. \\
5 & Visibility, uniform disk, binary stars, SII SNR, missing phase & Axes carry units; use \(B_\perp\) for the baseline, not the physical array separation. \\
6 & Detector response, jitter, dead time, background dilution, data rate & Normalize the response kernel; make the time units explicit. \\
7 & Pair count, time-shift background, covariance, and Fisher scaling & Draw accidental coincidences and physical correlations separately. \\
8 & Rayleigh curse, SPADE mode probabilities, QFI, and SII Fisher & In the small-separation limit, do not use a log axis to manufacture a false improvement. \\
9--13 & Radiation mechanisms, stars, compact objects, black holes, and transient toy models & Typical parameters come from the text; do not merely plot arbitrarily normalized curves. \\
14--17 & Propagation, polarization, new physics, CMB, and quantum-network resources & Ordinary astrophysical terms and new-physics terms must be drawn separately. \\
18--22 & Error budgets, science cases, teaching experiments, pitfalls, and roadmap & Captions should state the feasibility, the failure modes, or the decision boundary. \\
\bottomrule
\end{longtable}

\section{Minimum Requirements for Each Figure}
\label{appsec:e-figure-requirements}

\begin{longtable}{p{0.24\linewidth}p{0.66\linewidth}}
\toprule
Check item & Requirement \\
\midrule
Physical quantity & Axes state the quantity name and units clearly; dimensionless quantities must be explicitly labeled dimensionless or have units omitted. \\
Order of magnitude & At least one curve, point, or annotation should correspond to a real order of magnitude from the text. Pedagogically amplified signals must be flagged in the caption. \\
Symbols & Symbols in the figure should match the text, for example \(B_\perp\), \(\lambda\), \(\theta\), \(R_\gamma\), \(\Delta t\), \(\tau_c\). \\
Output & All main-text figures use PDF; no screenshots or low-resolution bitmaps. \\
Style & Use the unified style; avoid overly dense grids, too many colors, and decoration with no physical meaning. \\
Reproducibility & A comment at the top of the script or on the function should state the input parameters, default values, and output files. \\
\bottomrule
\end{longtable}

\section{Debugging Checklist}
\label{appsec:e-debug}

If a figure looks inconsistent with the text, check these issues first:

\begin{enumerate}[leftmargin=*]
\item Are units mixed up: mas, rad, m, nm, Hz, s are the easiest to get wrong.
\item Is the angular diameter being treated as an angular radius, or vice versa.
\item Is the baseline the projected baseline \(B_\perp\), rather than the physical array separation.
\item Is the spectral bandwidth \(\Delta\lambda\) being substituted directly for \(\Delta\nu\).
\item Are the correlation bin, the detector response, and the coherence time all in the same time unit.
\item Are the random numbers using a fixed seed; do the error bars of simulated figures change substantially with the seed.
\item Do the generated file names match the \texttt{\string\includegraphics} paths in the LaTeX.
\end{enumerate}

\section{Minimum Submission Package for the Course Project}
\label{appsec:e-project-package}

When submitting the course project in Chapter~\ref{chap:e26}, include at least five files or objects: an event-table description, the Python script, the generated PDF figures, a short report, and the null-test output. In the short report, every figure must state its input parameters, units, physical meaning, and failure modes; pasting the figure alone is not enough.
